%==============================
%	Importation des packages
%==============================
% --- Encodage & langue ---
\usepackage[utf8]{inputenc}
\usepackage[T1]{fontenc}
\usepackage[french]{babel}

% --- Maths ---
\usepackage{amsmath, amssymb, amsthm, amsfonts}
\usepackage{bm}
\usepackage{stmaryrd}

% --- Mise en page ---
\usepackage[margin=1in]{geometry} %marge
\usepackage{lmodern} %police d'écriture vectoriel
\usepackage[colorlinks=true, linkcolor=AntiDAcolor, urlcolor=AntiDAcolor]{hyperref} %référence
\usepackage{graphicx} %images
\usepackage{todonotes} %lecture note
\usepackage{fancyhdr} %en-tête
\usepackage{titlesec} %format section
\usepackage{marginnote}
\usepackage{enumitem}
\usepackage{setspace} %interligne

\usepackage{xcolor} %couleur perso
\usepackage{tcolorbox}
\usepackage{mdframed}
\tcbuselibrary{skins, breakable,theorems}

%Définition des couleurs
\definecolor{DAcolor}{RGB}{31, 78, 95}
\definecolor{AntiDAcolor}{RGB}{232, 89, 12}


%===================================
%	Implémentation des lectures
%===================================
\setlength{\marginparwidth}{2cm}

\setuptodonotes{noline, color=DAcolor!20, size=\tiny}
\makeatletter
\@todonotes@SetTodoListName{Cours}
\makeatother

\newcounter{lecture}


%=================
%	Commandes
%=================
% Ensembles
\newcommand{\N}{\mathbb{N}}
\newcommand{\Z}{\mathbb{Z}}
\newcommand{\Q}{\mathbb{Q}}
\newcommand{\R}{\mathbb{R}}
\newcommand{\C}{\mathbb{C}}
\newcommand{\K}{\mathbb{K}}


\newcommand{\lecture}[1]{%
	\stepcounter{lecture}%
	\todo[caption={Cours n° \thelecture -- #1}]{%
		Cours n° \thelecture
	}%
}

\DeclareMathOperator{\Vol}{vol}
\newcommand{\et}{\quad \text{et} \quad}
\newcommand{\Rn}{\R^n}
\newcommand{\Rm}{\R^m}		

\newcommand{\norme}[1]{\| #1 \|}
\newcommand{\entiers}{\llbracket 1, n \rrbracket}

\newcommand{\om}[1]{m^*\left( #1 \right)
}

\newcommand{\prf}[1]{
	\begin{proof}
		#1
	\end{proof}
	\vspace{0.5cm}
}




%========================================
%	Implémentation des environnements
%========================================
\newtheoremstyle{coloredtitle}
{}
{}                          % espace avant/après
{\itshape}                    % corps en italique (comme plain)
{}                            % pas d'indentation
{\bfseries}                   % titre en gras (comme plain)
{.}                           % ponctuation après le titre
{ }                           % espace après
{\textcolor{DAcolor}{\bfseries\thmname{#1}\thmnumber{ #2}\thmnote{ (#3)}}}

          % Format du titre

\newtheoremstyle{coloredtitlenoint}
{}
{}                          % espace avant/après
{}                    % corps pas en italique (definition)
{}                            % pas d'indentation
{\bfseries}                   % titre en gras (comme plain)
{.}                           % ponctuation après le titre
{ }                           % espace après
{\textcolor{DAcolor}{\bfseries\thmname{#1}\thmnumber{ #2}\thmnote{ (#3)}}}


\theoremstyle{coloredtitle}
%\newtheorem{theorem}{Théorème}[chapter]
\newmdtheoremenv[
backgroundcolor=DAcolor!15,
fontcolor=black,
linecolor=DAcolor,
linewidth=0.5pt
]{theorem}{Théorème}[chapter]

\newtheorem{lemme}{Lemme}[chapter]
\newtheorem{proposition}{Proposition}[chapter]

\theoremstyle{coloredtitlenoint}
\newtheorem{definition}{Définition}[chapter]
\newtheorem{exemple}{Exemple}[chapter]
\newtheorem{remarque}{Remarque}[chapter]

\setcounter{secnumdepth}{3}

\tcolorboxenvironment{definition}{
	blanker, breakable,
	left=4mm,
	borderline west={2pt}{0pt}{DAcolor} 
}

\tcolorboxenvironment{theorem}{
	blanker, breakable,
	left=0.6mm,
	borderline west={4pt}{0pt}{DAcolor} 
}
\tcolorboxenvironment{proposition}{
	blanker, breakable,
	left=4mm,
	borderline west={2pt}{0pt}{DAcolor} 
}

%redéfinition de l'environnement "proof"
\renewenvironment{proof}
{\noindent \textsc{Démonstration} :\newline}
{\qed}


%==================
%	Mise en page et structure
%=====================================
\pagestyle{fancy}
\fancyhf{}  % efface tout

% Numéro de page en coin extérieur bas
\fancyfoot[LE]{\thepage}  % Left Even  → page de gauche
\fancyfoot[RO]{\thepage}  % Right Odd  → page de droite

% Optionnel — nom du cours en coin intérieur haut
%\fancyhead[RE, LO]{\small\textit{\matiere}}
\fancyhead[RO, LE]{\leftmark}
\renewcommand{\headrulewidth}{0pt}

%inter-ligne
\setstretch{1.1}

%configuration de l'indentation des listes
%\setlist{leftmargin=2\leftmargini}

%configuration de la numérotaion des équations
\numberwithin{equation}{section}

%\setlist[description]{font=\normalfont}

%configuration des sections et sous-sections
\titleformat{\section}[block]   % format "block"
{\centering\large\bfseries\color{DAcolor}}     % centrage + taille + gras
{\thesection}                 % numéro de la sous-section
{1em}                            % espace entre numéro et texte
{}  

\titleformat{\chapter}    % section = "Chapitre"
[block]                 % format en bloc
{\centering\large\color{DAcolor}}      % style : centré, grand (pas en gras)
{\MakeUppercase{Chapitre \thechapter}} % numéro en majuscules
{0pt}                   % pas d’espace entre numéro et texte
{\vspace{0.5ex}\\\bfseries\huge}  % saut de ligne + texte en énorme



\titleformat{\subsection}[block]
{\large\bfseries\color{DAcolor}}
{\thesubsection}
{1em}
{}

